\subsection{COMPLEX PROJECTIVE SPACES}

With the notation recalled in Chapter VI, § 3, no. 1, we make the following definition, analogous to the definition of real projective spaces:

DEFINITION 1. The projective space \(\mathbf{P}_{n}(\mathbf{C})\) , endowed with the topology which is the quotient of that of \(\mathbf{C}_{n+1}^{*}\) by the equivalence relation \(\Delta_{n}(\mathbf{C})\) , is called complex projective space of \(n\) dimensions.

The projective space \(\mathbf{P}_{1}(\mathbf{C})\) is called the complex projective line, and \(\mathbf{P}_{2}(\mathbf{C})\) is called the complex projective plane.

Most of the arguments relating to real projective spaces extend with very slight modifications to complex projective spaces.

In the first place we see that the topological space \(\mathbf{P}_{n}(\mathbf{C})\) is Hausdorff by the argument of Proposition 1 of Chapter VI, § 3, no. 1, which applies word for word simply by replacing R by C. Again, the proof of Proposition 2 of Chapter VI, § 3, no. 1 shows that \(\mathbf{P}_{n}(\mathbf{C})\) is compact and connected, and homeomorphic to the quotient of the sphere \(S_{2n+1}\) (considered as embedded in the space \(C_{n+1}^{*}\) , identified with \(R_{2n+2}^{*}\) ) by the equivalence relation induced on this sphere by \(\Delta_{n}(\mathbf{C})\) ; the only point of difference is that now, if \(n \geqslant 0\) , the equivalence classes for this relation are homeomorphic to the circle \(S_{1}\) .

For this reason, Proposition 3 of Chapter VI, § 3, no. 1 has no analogue for complex projective spaces.

Next one shows, as in Chapter VI, § 3, no. 2, that every \(p\) -dimensional projective linear variety in the space \(\mathbf{P}_n(\mathbf{C})\) is a closed set, homeomorphic to \(\mathbf{P}_p(\mathbf{C})\) , and that its complement is dense in \(\mathbf{P}_n(\mathbf{C})\) if \(p < n\) . The proof of Proposition 5 of Chapter VI, § 3, no. 2 can be transposed as it stands, simply by substituting \(\mathbf{C}\) for \(\mathbf{R}\) , and shows that (if \(n \geqslant 0\) ) the complement of a projective hyperplane in \(\mathbf{P}_n(\mathbf{C})\) is homeomorphic to \(\mathbf{C}^n\) , and therefore that every point has a neighbourhood homeomorphic to \(\mathbf{C}^n\) . This result allows us to embed complex number space \(\mathbf{C}^n\) in complex projective space \(\mathbf{P}_n(\mathbf{C})\) , by identifying \(\mathbf{C}^n\) with the complement of a projective hyperplane, called the ``hyperplane at infinity'' (usually the hyperplane whose equation is \(x_0 = 0\) ). In the particular case \(n = 1\) , the hyperplane at infinity is a point, and Alexandroff's theorem shows that \(\mathbf{P}_1(\mathbf{C})\) is homeomorphic to the space \(\tilde{\mathbf{C}}\) obtained by compactifying the locally compact space \(\mathbf{C}\) by adjoining a ``point at infinity'', denoted by \(\infty\) . Proposition 4 of Chapter VI, § 2, no. 4 then shows that the complex projective line \(\mathbf{P}_1(\mathbf{C})\) is homeomorphic to the sphere \(S_2\) .

We leave to the reader the task of enunciating the results analogous to those of Chapter VI, § 3, no. 4, for functions which take their values in C.

Consider the space \(R^{n+1}\) as embedded in \(C^{n+1}\) (no. 1). Let f be the canonical mapping of \(C_{n+1}^{*}\) onto its quotient space \(\mathbf{P}_{n}(\mathbf{C})\) . The subspace \(f(\mathbf{R}_{n+1}^{*})\) consists of the points of \(\mathbf{P}_{n}(\mathbf{C})\) which have at least one system of real homogeneous coordinates; let us show that \(f(\mathbf{R}_{n+1}^{*})\) is homeomorphic to real projective space \(\mathbf{P}_{n}(\mathbf{R})\) , which will allow us to consider the space \(\mathbf{P}_{n}(\mathbf{R})\) as embedded in \(\mathbf{P}_{n}(\mathbf{C})\) . Now the relation induced by \(\Delta_{n}(\mathbf{C})\) on \(R_{n+1}^{*}\) is \(\Delta_{n}(\mathbf{R})\) ; the canonical mapping \(\varphi\) of

\[
\mathbf {R} _ {n + 1} ^ {*} / \Delta_ {n} (\mathbf {R}) = \mathbf {P} _ {n} (\mathbf {R})
\]

onto \(f(\mathbf{R}_{n+1}^*)\) is continuous (Chapter I, § 3, no. 6, Proposition 10); since

\(P_{n}(\mathbf{R})\) is compact, \(\varphi\) must be a homeomorphism (Chapter I, § 9, no. 4, Theorem 2, Corollary 2).

We can also prove that \(\varphi\) is bicontinuous without using the compactness of \(\mathbf{P}_n(\mathbf{R})\) , by invoking the criterion of Proposition 10 of Chapter I, § 3, no. 6 (see Exercise 3).

Since every vector subspace of \(p + 1\) dimensions of \(R^{n+1}\) generates a complex vector subspace of \(p + 1\) dimensions in \(C^{n+1}\) , we see that every projective linear variety V of p dimensions in \(\mathbf{P}_{n}(\mathbf{R})\) (V is called a real projective linear variety) generates a projective linear variety \(V'\) of p dimensions in \(\mathbf{P}_{n}(\mathbf{C})\) ( \(V'\) is called a complex projective linear variety), such that V is the trace of \(V'\) on \(\mathbf{P}_{n}(\mathbf{R})\) . Moreover, every system of (homogeneous) equations of V is also a system of (homogeneous) equations of \(V'\) when we allow the variables to take complex values.

